\section{Related work}\label{sec:related}

\subsection{Efficient and sparse attention}

The literature on cheaper attention is organised by pattern
\cite{tay2022efficient}: windowed attention, global tokens
\cite{beltagy2020longformer,zaheer2020bigbird}, low-rank and kernel
approximations, and learned selection \cite{zhou2021informer}. The usual
argument is a complexity statement at the level of the score matrix. Longformer
combines a sliding window with a small task-specific set of global tokens and
argues linear cost in $n$; Big Bird adds random links and proves that the
resulting pattern retains the expressive power of full attention, again with a
number of global tokens that does not grow with the sequence. Informer selects a
fixed number of queries and keys per step, so its budget per query is constant
in $n$. These analyses are correct for the designs they describe.

What is missing is a statement that separates the designs. ``Sparse'' is used
for patterns whose global-key budget is constant and for patterns whose budget
grows with $n$, and the two have different complexity classes with the same
name. Section~\ref{sec:theory} supplies the missing parameter and shows that it
is the budget, not the pattern, that determines the class. A second gap is the
distance between the mask and the kernel: complexity is stated for the set of
attended entries, whereas hardware executes tiles. The gap is not a detail; at
the settings we measure it changes the executed density by up to a factor
\BlockInflationGranMax (Section~\ref{sec:granularity}).

The strongest dense baseline is also not the one usually reported. IO-aware
exact attention \cite{dao2022flashattention} avoids materialising the score
matrix, so comparisons against a naive dense implementation overstate the saving
of a sparse pattern by the whole memory of that matrix. We therefore report both
dense baselines and treat the fused kernel as the primary competitor
(Section~\ref{sec:device}). The same reasoning applies to the
memory-minimal-but-launch-heavy variant of our own sparse kernel, which we
measure separately rather than choosing whichever is convenient.

\subsection{Attention as an object of analysis}

A complementary line studies attention as a computational primitive with its own
capacity and structure, for instance the analysis of attention building blocks
as associative memories with measurable storage limits
\cite{baldi2023quarks}. That work asks what a given number of heads and
dimensions can represent. We ask the dual question of what a given mask costs to
execute, and we connect the two through a retained-mass bound
(Theorem~\ref{thm:mass}): the price of a cheaper mask is stated in the same
currency, namely the probability mass it keeps, as the capacity question. The two
answers are not substitutes; a pattern can be expressive and unaffordable, or
cheap and lossy, and a design decision needs both numbers.

\subsection{Cost-sensitive learning and tabular decisions}

Threshold choice under asymmetric costs is classical
\cite{elkan2001foundations}: for known costs the optimal decision is a
likelihood-ratio threshold, and in practice the threshold is estimated on
validation data. We adopt that framework unchanged and use it as the evaluation
lens, because it is the only lens under which an efficiency question is a
decision question. Our contribution here is not a new cost-sensitive rule but
the observation that a substitution which improves a ranking metric can still
raise \eqref{eq:cost}, and that the two must be reported together.

On the modelling side, gradient-boosted trees remain the strongest family on
medium-sized tabular problems \cite{chen2016xgboost,ke2017lightgbm}, and careful
benchmarks show that the advantage survives modern architectures and tuning
\cite{gorishniy2021revisiting,grinsztajn2022why}. Our results are consistent
with that literature and extend it in a specific direction: we ask whether a
\emph{pretrained} model, adapted with a modern parameter-efficient method, closes
the gap when supervision is scarce, and we separate the two things adaptation can
repair: the output interface and the decision boundary.

\subsection{Language models on structured and tabular data}

Serialising a row as text and prompting a language model is now a standard
few-shot baseline for tabular classification \cite{hegselmann2023tabllm}. Two
measurement problems recur in that literature and we address both. First, a
generated answer must be parsed, and parse failures are frequently imputed
rather than counted, which manufactures discrimination from a constant. We count
them, and we show that a parse-free logit interface removes the failure mode at a
fraction of the cost. Second, adaptation and discrimination are usually reported
as one number; we report them as two, because they move in opposite directions in
our experiment.

\subsection{Verifiability and the operational context}

The system that motivated this study also had to convince an auditor that the
score it received was the score a committed model produced. Hardware-rooted
attestation \cite{costan2016sgx} and proof-of-learning
\cite{jia2021proof} are the two standard answers. Neither is measurable on our
device, so we report the commitment-and-challenge protocol we \emph{can} measure
and mark the rest as blocked (Section~\ref{sec:system}). We mention the
operational literature on hidden debt in learned systems \cite{sculley2015hidden}
because it explains why the defaults we study are attractive: each removes a
component that is hard to operate, and the cost of the removal is paid later and
elsewhere.

\subsection{Reproducibility as a design constraint}

Reported empirical AI results are unevenly reproducible, and structured
programmes improve this measurably \cite{pineau2021improving}. We take the
stronger position that a result which cannot refute its own author is not yet a
result, and we build the apparatus accordingly: claims carry status, refuted
hypotheses are recorded as refuted rather than deleted, and blocked measurements
carry the command that would unblock them. Section~\ref{sec:audit} describes the
mechanism and Section~\ref{sec:attn_withdraw} shows it firing on our own work.

\subsection{The closest prior work, and how this paper differs}

The three works closest to ours are the survey of efficient transformers
\cite{tay2022efficient}, the Big Bird analysis of sparse attention with global
tokens \cite{zaheer2020bigbird}, and the cost-sensitive learning framework of
\cite{elkan2001foundations}.

Relative to \cite{tay2022efficient}: that survey provides the taxonomy we use and
states complexity per pattern, at the level of attended entries and for the
designs as published. We add the parameter the taxonomy does not record, the
number of keys visible to every query as a function of $n$, and show that it
decides the complexity class, not the name of the pattern; we then measure
the executed rather than the nominal cost, which introduces a second factor, the
block-granularity inflation of Proposition~\ref{prop:block}.

Relative to \cite{zaheer2020bigbird}: that paper establishes that a combination
of window, global and random links preserves expressive power at linear cost, and
its global set is of fixed size. Our result is the complement: keeping the
window and the sinks but replacing the fixed global set with a \emph{periodic}
one, a natural choice when no task-specific global set is known, moves the
pattern out of the linear class, and we quantify the price on a real device. The
two results are consistent; together they say that the choice of how a position
becomes globally visible is a complexity decision, not a stylistic one.

Relative to \cite{elkan2001foundations}: the cost-sensitive framework is used
unchanged. What is new is the object it is applied to: we ask whether a
pretrained language model can replace the scorer under that criterion when
supervision is scarce, and we answer with a dissociation between interface
repair and discrimination loss that the framework alone does not predict.

Finally, relative to the tabular benchmarking literature
\cite{gorishniy2021revisiting,grinsztajn2022why}, our setting differs in the
budget rather than in the model zoo: we deliberately weaken the supervised
baseline (a fixed feature and row budget) so that the comparison favours the
pretrained model, and the gap nevertheless remains an order of magnitude in
AUC-PR. That is the sense in which we read our result as a supervision bound
rather than as another entry in a leaderboard.
